\chapter{Construção do sistema: passo a passo}

Este capítulo relata, em ordem cronológica, os vinte incrementos que levaram o Robofael da primeira
versão funcional à versão 3.8.0. Para cada incremento são descritos o pedido que o motivou, o que foi
construído, as decisões tomadas, os problemas encontrados e a forma de verificação. A
Tabela~\ref{tab:evolucao} resume a evolução.

\begin{singlespace}
\legendatabela{Evolução do sistema por incremento}\label{tab:evolucao}
\small
\begin{tabularx}{\textwidth}{C{0.8cm}C{1.6cm}LC{2.2cm}}
\toprule
\textbf{Nº} & \textbf{Versão} & \textbf{Incremento} & \textbf{Testes} \\ \midrule
\endfirsthead
\toprule
\textbf{Nº} & \textbf{Versão} & \textbf{Incremento} & \textbf{Testes} \\ \midrule
\endhead
1  & --     & Protótipo: avatar a partir das fotos, voz, agenda e modelo de linguagem & -- \\
2  & --     & Busto centralizado e boca suave & -- \\
3  & --     & Perfil de voz (ajuste de tom e ritmo) & -- \\
4  & --     & Nitidez, camisa futurista e boca natural & -- \\
5  & 1.0.0  & Estrutura profissional, documentação e testes unitários & 98 \\
6  & 2.0.0  & Refatoração completa: arquitetura hexagonal e pipeline em fluxo contínuo & 168 \\
7  & 2.1.0  & Renomeação para Robofael & 168 \\
8  & --     & Gravação da voz do usuário e criação do perfil real & 168 \\
9  & 2.2.0  & Ativação por voz local (“ok Robofael”) & 211 \\
10 & 2.3.0  & Nova frase de ativação (“ok Fael”) com gramática restrita & 247 \\
11 & 3.0.0  & Memória de longo prazo, “o que eu fiz”, janela fixa, configurações e programas & 351 \\
12 & 3.0.1  & Ganho automático do microfone e aviso de voz baixa & 357 \\
13 & 3.1.0  & Internacionalização (português e inglês) & 508 \\
14 & 3.2.0  & Instalador padrão para Ubuntu & 508 \\
15 & 3.3.0  & Revisão SOLID e código limpo, com testes de arquitetura & 529 \\
16 & 3.4.0  & Desempenho e fluidez guiados por medição; tipagem completa & 546 \\
17 & 3.5.0  & Nome técnico \cod{robofael}, com compatibilidade com os nomes antigos & 549 \\
18 & 3.6.0  & Modelo de linguagem local (Ollama) com respostas em camadas & 576 \\
19 & 3.7.0  & Escolha do microfone e abertura junto com o computador & 601 \\
20 & 3.8.0  & Local e backup do banco; modelo em uso e modelos compatíveis & 643 \\
\bottomrule
\end{tabularx}
\fonte{elaboração própria (2026), a partir do registro de mudanças do projeto.}
\end{singlespace}

\section{Incremento 1: protótipo funcional}

\textbf{Pedido.} Criar um programa em Python que manipulasse as imagens do usuário, presentes na pasta
do projeto, para transformá-las no avatar de um assistente pessoal de atividades diárias que ouve e
responde por voz.

\textbf{Construção.} Foram usadas duas fotografias: a do crachá e uma \en{selfie}, ambas já sem
fundo. Em cada uma foram medidas manualmente as posições da boca e do queixo. A animação foi obtida pela
deformação suave da região da mandíbula para baixo, com o preenchimento do espaço aberto por um interior
escuro de boca, em oito níveis de abertura. Um brilho ciano nas bordas deu ao avatar a aparência de
assistente digital. A voz foi implementada com reconhecimento de fala em português pelo serviço do Google
e síntese pela voz neural \cod{pt-BR-AntonioNeural}, com o gTTS como reserva. A abertura da boca passou a
acompanhar o volume real do áudio, quadro a quadro. A agenda de tarefas e lembretes falados era salva em
um arquivo JSON, e o modelo Claude, quando disponível, operava a agenda por meio de ferramentas; sem
chave de acesso, um modo offline atendia a comandos básicos.

\textbf{Verificação.} A interface foi executada com os dois avatares e a fala sincronizada com a
animação. A integração com o modelo de linguagem não pôde ser exercitada de imediato, por falta de
credenciais na máquina, e foi registrada como pendente.

\section{Incremento 2: busto centralizado e boca suave}

\textbf{Pedido.} Exibir apenas um busto, centralizado, e suavizar a animação da boca.

\textbf{Construção.} Mediu-se o topo e o centro da cabeça em cada foto, e o recorte passou a ir de um
pouco acima da cabeça até os ombros, centralizado na cabeça, e não na foto inteira. Com isso, a
\en{selfie}, que estava deslocada para o lado, ficou centralizada. Uma moldura oval com bordas suaves
passou a esconder os cortes das fotos. Para a boca, o número de aberturas subiu de 8 para 16, mais
próximas entre si quando a boca está quase fechada, onde os saltos eram mais visíveis. Cada quadro passou
a ser desenhado no dobro do tamanho e depois reduzido (superamostragem), eliminando o serrilhado. O
volume da voz passou a ser lido 60 vezes por segundo, sem picos isolados nem ruído de fundo. A boca
passou a abrir rápido e fechar devagar, como na fala real, e a animação foi para cerca de 60 quadros por
segundo.

\textbf{Verificação.} Capturas de tela confirmaram o enquadramento e a passagem por aberturas
intermediárias durante a fala.

\section{Incremento 3: perfil de voz}

\textbf{Pedido.} Registrar a voz do usuário e fazer o assistente falar com o mesmo tom. Diante das
alternativas, ajustar apenas o tom de uma voz neural ou clonar o timbre localmente, o usuário escolheu
a primeira.

\textbf{Construção.} Foi criado um gravador de voz com interface própria, que apresenta frases para
leitura, detecta o fim de cada uma e permite regravar. Com pelo menos três gravações, o programa estima o
tom médio por autocorrelação e a velocidade de fala. Em seguida, escolhe a voz neural em português de tom
mais próximo e ajusta tom e ritmo até que a voz sintética fique a menos de 3 Hz da voz do usuário. O
perfil é salvo em arquivo JSON.

\textbf{Verificação.} O medidor de tom acertou sinais de teste com erro máximo de 1 Hz. Em uma
simulação com um “usuário” 25 Hz mais grave e 15\% mais lento, o perfil resultou em um ajuste de
$-22$ Hz, com o tom final a 2 Hz do alvo, e em um ritmo de $-14\%$.

\section{Incremento 4: nitidez, camisa futurista e boca natural}

\textbf{Pedido.} Aumentar a nitidez do avatar, trocar a roupa por uma camisa de corte futurista, tornar
o movimento da boca mais natural e manter o busto centralizado.

\textbf{Construção.} Constatou-se que a \en{selfie} sem fundo tinha apenas 433 por 577 pixels, mas o
arquivo original tinha 1200 por 1600 pixels na mesma proporção. Passou-se então a usar a imagem sem fundo
apenas como máscara de contorno, aplicada sobre a foto em alta resolução. A foto do crachá, disponível
apenas em 500 pixels, recebeu ampliação de alta qualidade, filtro de nitidez e leve aumento de
contraste. A camisa original foi substituída por uma camisa azul-grafite desenhada por cima, com gola
alta metálica de borda luminosa, costuras em V com brilho ciano, placas nos ombros com pequenos LEDs e
um emblema no peito. As dobras do tecido foram preservadas aproveitando-se a luz e a sombra da camisa
original. A boca deixou de ser um corte que descia e passou a ser uma deformação contínua: o lábio
superior sobe levemente, a mandíbula desce, os cantos ficam presos e a abertura tem formato de lente,
com sombra sob os dentes e língua visível nas aberturas maiores. Foram introduzidos três formatos de
vogal (arredondado para “o” e “u”, neutro para “a” e mais largo para “e” e “i”), escolhidos a partir do
timbre do áudio.

\textbf{Verificação.} Prévias geradas com as mesmas imagens e posições da interface, além da execução
da interface com fala simulada, sem erros.

\section{Incremento 5: estrutura profissional e testes (versão 1.0.0)}

\textbf{Pedido.} Organizar o projeto com estrutura profissional, pronta para versionamento, com um
arquivo README que explicasse a arquitetura, a forma de execução e os fluxos, e com testes unitários.

\textbf{Construção.} Os scripts soltos foram reorganizados em um pacote Python instalável, com
subpacotes para lógica (agenda, conversa, ferramentas, modo offline), áudio, avatar e interface, e com a
configuração centralizada e ajustável por variáveis de ambiente. As fotos foram movidas para uma pasta de
recursos com nomes legíveis. A lógica foi separada do microfone, do alto-falante e da tela, e o cliente
do modelo, o relógio e as pastas tornaram-se substituíveis por versões simuladas, o que viabilizou os
testes sem hardware. Os dados pessoais (agenda, perfil e gravações) e as imagens geradas foram colocados
em pastas excluídas do controle de versão. Foram criados um Makefile, o arquivo de projeto
\cod{pyproject.toml} e o README.

\textbf{Verificação.} 98 testes unitários aprovados em cerca de três segundos e análise de estilo sem
ocorrências.

\section{Incremento 6: refatoração completa (versão 2.0.0)}

\textbf{Pedido.} Refatorar o projeto do início ao fim, com nível de engenharia de grandes empresas de
tecnologia, aperfeiçoando a usabilidade, o desempenho e a organização arquitetural.

\textbf{Diagnóstico.} Antes de qualquer alteração, foi feito um diagnóstico. O fluxo era totalmente
serial: esperava-se a resposta inteira do modelo, sintetizava-se todo o áudio e só então a fala começava,
com três a seis segundos de silêncio. O microfone usava uma função da biblioteca que funcionava como
caixa-preta, sem nível de áudio nem cancelamento. Os estados eram cadeias de caracteres alteradas por
várias threads. A renderização do avatar bloqueava a abertura da janela. Os lembretes eram verificados
a cada 15 segundos. A escrita direta do JSON podia corromper o arquivo. Por fim, a classe principal da
interface conhecia todas as tecnologias.

\textbf{Construção.} A solução foi reescrita em arquitetura hexagonal, com:
\begin{enumerate}[label=\alph*)]
  \item domínio tipado com classes de dados no lugar de dicionários;
  \item máquina de estados explícita e barramento de eventos;
  \item pipeline em fluxo contínuo de três estágios, com divisor incremental de frases;
  \item áudio decodificado em memória e envelope da boca calculado antes da reprodução;
  \item cache em disco das falas sintetizadas;
  \item renderização paralela do avatar fora da thread da interface, com cache invalidado pelo resumo do
        código;
  \item inicialização paralela dos serviços em segundo plano;
  \item agendador que dorme até o próximo lembrete;
  \item detecção de fala própria, com medidor de nível real e pausa de fim de fala de 0,7 s;
  \item escrita atômica da agenda e preservação de arquivos corrompidos;
  \item histórico da conversa somente de acréscimo, com descarte por rodada;
  \item validação local das entradas das ferramentas.
\end{enumerate}
Foram adotados o verificador de tipos, a análise estática, os ganchos de \en{pre-commit} e a integração
contínua com servidor gráfico virtual.

\textbf{Problemas corrigidos.} Foram corrigidos cinco problemas: um cancelamento que se perdia quando o
ESPAÇO era pressionado duas vezes rapidamente; uma notificação perdida no agendador; um erro de ponto
flutuante no limiar de tempo da detecção de fala; o acesso à interface do gravador fora da thread
principal; e valores padrão da configuração incompatíveis com classes de dados compactas.

\textbf{Verificação.} 168 testes, incluindo um teste que prova que a primeira frase toca antes de o
modelo terminar de gerar a resposta. A janela passou a aparecer em 0,8 s, e a primeira palavra no modo
offline, em cerca de 1,3 s.

\section{Incremento 7: renomeação para Robofael (versão 2.1.0)}

\textbf{Pedido.} Mudar o nome do assistente de EuRobô para Robofael.

\textbf{Construção.} O nome passou a ficar centralizado em um único campo da configuração, ajustável
por variável de ambiente, e foi aplicado ao título da janela, à saudação, às instruções do modelo, à
ferramenta de gravação, aos registros e à documentação. Optou-se por manter o nome técnico do pacote
(\cod{eurobo}) para não quebrar importações, comandos e caminhos existentes. A renomeação técnica
para \cod{robofael} só ocorreu na versão 3.5.0 (Incremento 17), mantendo os comandos e as variáveis
de ambiente antigos como alternativas.

\textbf{Verificação.} Testes do nome padrão, da troca por variável e da apresentação do modelo como
Robofael.

\section{Incremento 8: gravação da voz real do usuário}

\textbf{Pedido.} Gravar a voz do usuário.

\textbf{Construção e resultado.} A ferramenta de gravação foi aberta, e o usuário gravou as cinco
frases. O tom médio medido foi de 116 Hz. A voz base mais próxima foi a \cod{pt-BR-AntonioNeural}, com
137 Hz. O ajuste de tom de $-21$ Hz levou a voz sintética a 118 Hz, a 2 Hz da voz do usuário, e o ritmo
ficou em $-26\%$. As gravações tiveram entre 3,9 e 5,2 segundos de fala cada uma, com tom entre 95 e
119 Hz, e o perfil usa a mediana das cinco para reduzir o efeito de frases interrogativas.

\textbf{Problemas e diagnóstico.} Durante a primeira execução real, o modelo de linguagem não
conectou. A investigação mostrou que a configuração do terminal do usuário direcionava as chamadas da API
para um servidor local de outro fornecedor, onde o modelo não existia. A mensagem de aviso foi
aprimorada para indicar essa causa, sem alterar a configuração pessoal do usuário. Também foi eliminado
o ruído de mensagens nativas do servidor de áudio JACK que aparecia no terminal a cada abertura do
microfone.

\section{Incremento 9: ativação por voz (versão 2.2.0)}

\textbf{Pedido.} Acionar a captura de voz com a frase “ok Robofael”, em vez de pressionar o ESPAÇO.
Entre o reconhecimento local e o remoto da frase de ativação, o usuário escolheu o local.

\textbf{Construção.} Foi implementada a escuta passiva local com o Vosk e um modelo compacto de
31 MB para português, para que o áudio do ambiente não deixasse o computador. Como o modelo não conhece
a palavra “Robofael” e a transcreve por aproximação (por exemplo, “o que roubou fel”), a detecção foi
feita por semelhança fonética, calibrada com transcrições reais: ativa com a frase sozinha, com a frase
seguida de comando e com a frase dita após outra conversa, mas não com “ok Rafael” nem com conversas
comuns. Quando o reconhecedor remoto erra a frase de ativação, como em “Ok Google Fael, me lembra…”, o
comando é recuperado pelo alinhamento das duas transcrições. Um sinal sonoro de dois tons confirma a
ativação.

\textbf{Problemas corrigidos.} A biblioteca de áudio capturava os sinais de interrupção do sistema, e
o programa não fechava com Ctrl+C nem com SIGTERM. A solução foi desativar os tratadores de sinal da
biblioteca e tratar os sinais no laço da interface.

\textbf{Verificação.} 211 testes; reconhecimento real exercitado com a voz neural do perfil injetada no
lugar do microfone.

\section{Incremento 10: frase “ok Fael” e gramática restrita (versão 2.3.0)}

\textbf{Pedido.} Trocar a frase de ativação para “ok Fael”.

\textbf{Construção.} A palavra “Fael”, curta e fora do vocabulário do modelo local, era simplesmente
descartada: “ok Fael” virava “ok” ou “o que foi”. A transcrição livre foi então substituída por uma
gramática restrita, na qual o reconhecedor só pode escolher entre variantes aceitas de “ok Fael” e
frases distratoras parecidas e frequentes. Com isso, o que não é a frase de ativação é absorvido pelas
distratoras. A frase de ativação passou a ser configurável e independente do nome do assistente. Como a
pausa da vírgula dividia a fala em duas partes, a escuta passou a continuar gravando enquanto houvesse
voz após a frase. O comando também passou a ser limpo da frase quando o reconhecedor remoto a transcreve
de forma inesperada, como em “o que Israel anota tarefa…”, que resulta em “anota tarefa…”.

\textbf{Verificação.} Calibração com cinco vozes sintéticas: zero falsos positivos em 115 frases
parecidas e 8 acertos em 8 tentativas na voz do perfil do usuário; 247 testes.

\section{Incremento 11: memória de longo prazo e outras funções (versão 3.0.0)}

\textbf{Pedido.} Adicionar um banco SQLite como memória de longo prazo, com histórico auditável de
tarefas e lembretes, para que o assistente respondesse o que o usuário fez em uma data ou semana do
passado. O pedido incluía ainda tornar fixo o tamanho da janela, porque a tela cheia quebrava a
renderização; criar um botão de configuração de avatar e voz; permitir abrir programas instalados no
Linux; e garantir a integridade de todos os testes.

\textbf{Construção.} O pedido foi dividido em seis frentes:
\begin{enumerate}[label=\alph*)]
  \item \textbf{memória em SQLite}: agenda, histórico e preferências em um único arquivo; histórico
        somente de acréscimo, com gatilhos que recusam alteração e remoção e com cadeia de resumos SHA-256;
        estado e histórico gravados na mesma transação; esquema versionado; importação automática do
        arquivo JSON anterior, renomeado e preservado; verificação de integridade por comando;
  \item \textbf{“o que eu fiz”}: interpretação de períodos (dezenove formas testadas) e cálculo do
        resultado líquido no período, disponível no modo offline e como ferramenta do modelo;
  \item \textbf{janela de tamanho fixo}: tamanhos mínimo e máximo iguais e redimensionamento desativado;
  \item \textbf{configurações}: diálogo com avatar, voz, tom, ritmo, botão para ouvir, uso do perfil
        gravado e ativação por voz, com cancelamento que desfaz o teste e preferências persistidas;
  \item \textbf{programas}: leitura das entradas \cod{.desktop}, sinônimos, navegador padrão,
        desambiguação e abertura sem interpretador de comandos;
  \item \textbf{testes}: cobertura de todas as frentes, incluindo a detecção de quatro tipos de
        adulteração do histórico.
\end{enumerate}

\textbf{Problemas corrigidos.} Um teste revelou que “anota tarefa abrir conta no banco” abria um
programa, em vez de criar a tarefa; os comandos de agenda passaram a ter prioridade. A análise estática
apontou uma função anônima que usava uma variável já descartada pelo tratamento de exceção. As
variáveis da interface eram destruídas fora da thread principal pelo coletor de lixo e passaram a ser
liberadas explicitamente.

\textbf{Verificação.} 351 testes aprovados em três execuções seguidas e execução real com uma cópia dos
dados do usuário, na qual a migração, a consulta ao histórico e a verificação de integridade
funcionaram.

\section{Incremento 12: ganho automático (versão 3.0.1)}

\textbf{Contexto.} Os primeiros testes com a voz real do usuário mostraram que a frase de ativação não
era reconhecida: o reconhecedor ouvia apenas “ok” e “que”, e o pico do sinal ficava em cerca de 580,
quando o necessário era superior a 1000. A investigação, feita somente por leitura, revelou que o
microfone em uso era o interno do computador, com entrada em 37\%, enquanto um microfone USB conectado
não estava sendo usado.

\textbf{Construção.} Foi implementado um controle automático de ganho que amplifica até 8 vezes a voz
que chega baixa, resistente a picos isolados, e um aviso de voz baixa no terminal e na tela, com
orientação ao usuário. As configurações de áudio do sistema não foram alteradas; os comandos para
ajustá-las foram fornecidos ao usuário.

\textbf{Verificação.} 357 testes; após o ajuste, os registros mostraram nove ativações bem-sucedidas
com a voz real do usuário.

\section{Incremento 13: internacionalização (versão 3.1.0)}

\textbf{Pedido.} Permitir que o assistente responda em inglês quando configurado.

\textbf{Construção.} Todas as mensagens foram reunidas em um catálogo único, com as duas línguas lado a
lado, e um tradutor foi injetado em todos os componentes, com troca ao vivo. O idioma passou a afetar as
respostas, a interface, a voz neural (Antonio em português; Andrew, Brian, Guy ou Jenny em inglês), a
voz de reserva, o reconhecimento de fala e as instruções do modelo, cuja conversa recomeça na troca.
Os comandos passaram a ser entendidos nos dois idiomas, independentemente do idioma das respostas, e os
períodos em inglês foram suportados, com datas numéricas seguindo a convenção de cada idioma.

\textbf{Problemas corrigidos.} O diálogo de configurações podia falhar ao se tornar modal antes de
aparecer na tela e passou a tentar novamente. Além disso, as instruções ao modelo deixaram de atribuir
gênero ao usuário.

\textbf{Verificação.} 508 testes, inclusive um que garante que nenhuma tradução está faltando e que os
campos de formatação coincidem nas duas línguas; execução real em inglês com uma cópia dos dados.

\section{Incremento 14: instalador para Ubuntu (versão 3.2.0)}

\textbf{Pedido.} Criar um instalador padrão para ambientes Ubuntu.

\textbf{Construção.} Optou-se por um pacote \cod{.deb}, o formato nativo da distribuição
\cite{debian}. O pacote contém o código, as fotos, o modelo de ativação e todas as dependências Python
em formato \en{wheel}, o que permite instalar sem internet e de forma determinística. As dependências de
sistema, como Python, Tk, PyAudio e PortAudio, são declaradas para o gerenciador de pacotes. Na
instalação, um ambiente Python isolado é criado em \cod{/opt/robofael} e verificado. O pacote fornece
os comandos \cod{robofael} e \cod{robofael-gravar-voz}, atalhos no menu de aplicativos e um ícone
gerado a partir do próprio avatar. Os dados de cada usuário ficam em \cod{~/.local/share/robofael} e o
cache em \cod{~/.cache/robofael}, conforme as convenções XDG. A remoção apaga o programa, mas preserva
os dados do usuário. A janela passou a declarar a classe \cod{robofael}, para que o ambiente gráfico a
associe ao ícone do menu.

\textbf{Problemas corrigidos.} Uma das dependências não era distribuída como \en{wheel} e passou a ser
construída localmente. Além disso, o ambiente isolado reutilizava uma versão da biblioteca de imagens do
sistema sem suporte à interface gráfica, o que foi resolvido pela instalação forçada das versões
incluídas no pacote.

\textbf{Verificação.} Um roteiro de teste instala o pacote em uma pasta temporária e verifica a
integridade dos arquivos, os atalhos, a criação do ambiente sem internet, os lançadores, os caminhos de
dados, o funcionamento do núcleo e a remoção. A janela foi aberta a partir dessa instalação simulada.

\section{Incremento 15: revisão SOLID e código limpo (versão 3.3.0)}

\textbf{Pedido.} Validar o projeto de ponta a ponta, implementando e verificando os conceitos SOLID e de
código limpo, garantindo que os testes continuassem funcionando.

\textbf{Construção.} Esse incremento é detalhado no Capítulo 7. Em síntese, foram feitas as seguintes
alterações:
\begin{enumerate}[label=\alph*)]
  \item o orquestrador foi dividido em três classes;
  \item as políticas de sessão saíram da raiz de composição e da interface;
  \item três despachos condicionais tornaram-se tabelas extensíveis;
  \item contratos foram criados ou movidos para que todas as dependências apontassem para dentro;
  \item uma violação do princípio de Liskov foi eliminada;
  \item funções longas foram decompostas;
  \item regras de complexidade foram incorporadas à análise estática;
  \item testes de arquitetura passaram a verificar as regras automaticamente.
\end{enumerate}

\textbf{Problema corrigido.} Durante a revisão, um teste de interface passou a falhar de forma
intermitente apenas na suíte completa. A investigação, feita com a captura das pilhas de execução no
momento da falha, revelou que a thread de renderização do avatar ficava bloqueada no destrutor de uma
imagem Tk de um teste anterior: o coletor de lixo do Python, acionado naquela thread, chamava a
biblioteca gráfica fora da thread principal. O defeito era latente e havia sido apenas exposto pela
mudança no momento da coleta. A correção recolhe esses objetos na thread principal ao fim de cada teste
de interface, e a fixture duplicada foi unificada.

\textbf{Verificação.} 529 testes, 85\% de cobertura, nenhuma função com complexidade ciclomática acima
de 10 e imagens do avatar idênticas, pixel a pixel, às anteriores à refatoração.

\section{Incremento 16: desempenho e fluidez guiados por medição (versão 3.4.0)}

\textbf{Pedido.} Refatorar novamente o projeto de ponta a ponta, aperfeiçoando a usabilidade e o
desempenho, com uma organização ainda mais moderna e profissional.

\textbf{Diagnóstico.} Como as revisões anteriores já haviam tratado a estrutura, este incremento começou
por medições do aplicativo real, e não pela leitura do código. Com o assistente parado, depois da
saudação, o processo consumia 29\% de um núcleo de processamento. Um microbenchmark do palco mostrou que
cada quadro levava cerca de 13 ms, perto do limite de 16,6 ms de um quadro a 60 por segundo. Medições
isoladas de cada elemento revelaram a causa: qualquer alteração sob o halo (600 por 600 pixels com
transparência) ou sob o avatar obrigava a biblioteca gráfica a refazer, em software, a mistura de
transparência de toda a área, e girar um único arco custava 8 ms. Além disso, a montagem dos 40 halos
bloqueava a thread da interface por 0,75 s antes de a janela aparecer, e a troca de avatar convertia
36 imagens de uma vez, com um engasgo de cerca de 170 ms. Na fala, a medição do áudio sintetizado
mostrou que cada frase trazia cerca de 0,25 s de silêncio no início e de 0,9 a 1,2 s no fim, o que, na
reprodução frase a frase, se transformava em pausas artificiais.

\textbf{Construção.} Foram realizadas as seguintes alterações:
\begin{enumerate}[label=\alph*)]
  \item \textbf{canvas diferencial}: uma camada entre o palco e a biblioteca gráfica guarda o último valor
        enviado de cada item e só repassa coordenadas, em pixels inteiros, e opções que mudaram;
  \item \textbf{halo opaco}: o halo passou a ser gerado já misturado com a cor de fundo, de modo que a
        biblioteca copia os pixels em vez de misturar transparência; o redesenho sob o halo caiu de
        cerca de 8 ms para 0,1 ms;
  \item \textbf{taxa de quadros adaptativa}: um laço de quadros extraído da janela principal roda a 60
        quadros por segundo com o assistente ouvindo, pensando ou falando, ou durante transições de cor,
        e a 30 em repouso, agendando pelo relógio para não acumular atraso;
  \item \textbf{fila de trabalho com orçamento}: os halos passaram a ser gerados em segundo plano, e a
        conversão de halos e quadros do avatar para a biblioteca gráfica é feita aos poucos, com no
        máximo 4 ms por quadro;
  \item \textbf{corte do silêncio}: uma função de processamento de sinal localiza o trecho com fala e
        mantém margens de 30 ms antes e 150 ms depois, preservando consoantes fracas no início;
  \item \textbf{síntese paralela}: até duas frases são sintetizadas ao mesmo tempo e tocadas na ordem
        original; a decodificação continua em uma única thread, porque o mixer de áudio não é feito
        para uso concorrente;
  \item \textbf{pedido digitado durante a fala}: enviar um texto com o assistente ocupado passou a
        interromper o que ele fazia e atender o pedido novo, em vez de descartá-lo; se vários forem
        enviados, vale o mais recente, e uma trava garante que o pedido seja atendido exatamente uma vez;
  \item \textbf{tipagem completa}: todas as funções passaram a ter assinatura tipada, exigência agora
        verificada pelo mypy.
\end{enumerate}

\textbf{Problemas evitados.} Ao tipar a função de recorte do busto, uma primeira versão passou a usar
coordenadas arredondadas também no cálculo da posição da boca, o que a deslocaria em até meio pixel. A
comparação do resumo SHA-256 de todos os quadros, feita antes e depois, detectou a diferença, e a
correção restaurou as imagens idênticas pixel a pixel. Uma medição inicial do consumo em repouso também
se mostrou contaminada pela saudação falada e foi refeita depois que ela terminava, com a versão anterior
medida nas mesmas condições.

\textbf{Verificação.} 546 testes, incluindo uma regressão de desempenho determinística: um palco em
repouso, desenhado em um canvas falso que conta as chamadas, precisa enviar menos de 80 comandos por
segundo, e as imagens grandes só podem ser redesenhadas poucas vezes por segundo. A cobertura ficou em
85,9\%. No aplicativo real, o consumo em repouso caiu de 29\% para 5\% de um núcleo, a janela passou a
aparecer em 0,6 s (antes, de 1,7 a 2,4 s), a primeira palavra da saudação passou a sair de 0,4 a 0,8 s
mais cedo e a saudação completa ficou cerca de 1,6 s mais curta.

\section{Incremento 17: nome técnico robofael (versão 3.5.0)}

\textbf{Pedido.} Renomear o nome técnico do projeto, que ainda era \cod{eurobo}, para \cod{robofael},
sem quebrar comandos e variáveis que o usuário já usasse.

\textbf{Construção.} O pacote Python, os comandos, as variáveis de ambiente, o lançador do pacote
\cod{.deb} e a documentação passaram a usar o nome \cod{robofael}. Os nomes antigos foram mantidos como
alternativas: os comandos \cod{eurobo} e \cod{eurobo-gravar-voz} continuam instalados como apelidos, e
cada variável \cod{EUROBO_*} ainda é lida quando a \cod{ROBOFAEL_*} correspondente não está definida,
com um aviso de depreciação no registro. O lançador do pacote faz a mesma substituição, para que uma
pasta de dados já configurada pelo nome antigo continue sendo usada.

\textbf{Verificação.} 549 testes, incluindo a leitura dos nomes antigos com aviso e a precedência do
nome novo; um teste de arquitetura garante que o nome antigo só aparece no código de compatibilidade, e
o roteiro do pacote confere que o lançador respeita a variável antiga da pasta de dados.

\section{Incremento 18: modelo de linguagem local (versão 3.6.0)}

\textbf{Pedido.} Permitir que o assistente converse e entenda pedidos em linguagem livre sem o Claude,
com um modelo executado no próprio computador.

\textbf{Construção.} Foi criado um adaptador para o Ollama \cite{ollama}, servidor de modelos locais
acessado por HTTP, escrito apenas com a biblioteca padrão. A raiz de composição passou a escolher o
melhor cérebro disponível, na ordem Claude, modelo local e modo offline, com uma variável que restringe
a escolha. O modelo padrão, de 3 bilhões de parâmetros, foi escolhido por responder em tempo aceitável
no processador, sem placa de vídeo.

\textbf{Problema e decisão.} Os primeiros testes mostraram que um modelo pequeno é bom para conversar,
mas pouco confiável para escolher ferramentas: um pedido de curiosidade abriu o navegador e um
desabafo listou as tarefas. Por isso, a resposta passou a ser montada em camadas:
\begin{enumerate}[label=\alph*)]
  \item comandos que o modo offline reconhece são atendidos por ele, na hora e sempre do mesmo jeito;
  \item pedidos de agenda em linguagem livre vão ao modelo apenas com as ferramentas do assunto e
        somente as seguras (criar e consultar); concluir, remover e abrir programas ficam com os
        comandos explícitos;
  \item o restante é conversa e vai ao modelo sem ferramentas, com o texto em fluxo contínuo;
  \item se o servidor local não responder, a fala é atendida pelo modo offline, e o modelo só volta a
        ser tentado depois de um intervalo.
\end{enumerate}

Chamadas de ferramenta escritas como texto pelo modelo são executadas ou descartadas, mas nunca
faladas. O histórico é compactado a cada troca, e o modelo é carregado na memória antes da primeira
fala, com as mesmas opções da conversa, para evitar uma segunda carga.

\textbf{Verificação.} 576 testes. As camadas, as ferramentas seguras, a queda para o modo offline e o
cliente HTTP são testados contra um servidor falso; no uso real, a primeira palavra de uma conversa
saiu em 1 a 2,5 s, e pedidos com ferramentas levaram de 15 a 45 s no processador.

\section{Incremento 19: escolha do microfone e inicialização (versão 3.7.0)}

\textbf{Pedido.} Permitir selecionar o microfone entre os dispositivos do computador e configurar o
sistema para abrir junto com o computador ou pelo comando \cod{robofael} no terminal.

\textbf{Construção.} Duas portas novas foram criadas, uma para listar os dispositivos de áudio e outra
para a inicialização. Os microfones são listados pelo servidor de áudio (PipeWire \cite{pipewire} ou
PulseAudio), sem os monitores das saídas de som, e a escolha é salva nas preferências. Abrir o
dispositivo físico diretamente disputaria o aparelho com o servidor de áudio e exigiria taxa e canais
exatos; por isso, o microfone continua sendo aberto pelo dispositivo padrão, e a escolha é indicada ao
servidor por variáveis de ambiente definidas apenas durante a abertura, sob uma trava, para não serem
herdadas pelos programas que o assistente abre. A troca vale na próxima escuta: o microfone é
recalibrado e a escuta passiva é reaberta no novo dispositivo. A inicialização usa um atalho de sessão
no padrão de entradas \cod{.desktop} \cite{freedesktop}, com atraso de 5 s e verificação de que o
programa ainda existe, e atalhos em \cod{~/.local/bin} para o terminal. O estado é lido dos próprios
arquivos, e desligar só apaga atalhos criados para este programa. A tela de configurações passou a ter
duas colunas.

\textbf{Verificação.} 601 testes. Uma fixture passou a apontar a pasta pessoal para uma pasta
temporária, para que nenhum teste altere a configuração real do usuário. No aplicativo real, a lista
do servidor de áudio confirmou que a gravação passou a sair do microfone escolhido, e não do padrão.

\section{Incremento 20: banco, backup e modelos de linguagem (versão 3.8.0)}

\textbf{Pedido.} Permitir escolher onde o banco SQLite é salvo e fazer backup dele em qualquer local, e
exibir nas configurações qual modelo de linguagem está em uso e quais modelos compatíveis estão
disponíveis no Claude e no Ollama.

\textbf{Construção.} O banco ganhou as operações de mudar de lugar e de copiar, ambas com o assistente
em uso:
\begin{enumerate}[label=\alph*)]
  \item a cópia usa a interface de backup do SQLite \cite{sqlite}, consistente mesmo com gravações em
        andamento, e é conferida por uma verificação de integridade antes de ser aceita;
  \item ao mudar de lugar, a conexão é trocada sob a mesma trava das gravações, e o arquivo antigo é
        renomeado com a data, nunca apagado; a mudança é recusada se já houver um banco no destino;
  \item como o banco não pode guardar o próprio endereço, o local escolhido fica em um pequeno arquivo
        na pasta de dados; se o banco escolhido não estiver acessível na abertura, como em um disco
        desconectado, o assistente avisa e usa o banco padrão apenas naquela vez, sem criar um banco
        vazio no lugar;
  \item o backup é um arquivo único, sem o diário auxiliar, e informa quantos registros do histórico
        foram copiados.
\end{enumerate}

Para os modelos, cada provedor ganhou um catálogo: o do Claude lista os modelos liberados para a chave,
e o do Ollama lista os modelos instalados capazes de conversar e de chamar ferramentas, com o tamanho.
As consultas são feitas em paralelo, fora da thread da interface, e a falta de chave ou um servidor
fora do ar viram uma explicação na própria lista.

\textbf{Verificação.} 643 testes, incluindo a mudança do banco com gravações antes e depois, a recusa
de sobrescrever outro banco, o banco ausente na abertura, a verificação da cadeia de resumos na cópia e
os catálogos com respostas simuladas. No aplicativo real, a tela mostrou o modelo local em uso e os
dois modelos instalados, e o backup foi conferido.

\section{Lições aprendidas no processo}

O processo evidenciou algumas lições, discutidas novamente no Capítulo 9:

\begin{enumerate}[label=\alph*)]
  \item medições com dados reais, como o volume da voz do usuário e transcrições efetivamente produzidas
        pelos reconhecedores, foram decisivas para resolver problemas que testes com vozes sintéticas
        não revelavam;
  \item decisões com impacto na privacidade ou na experiência do usuário foram submetidas a ele, e
        configurações pessoais do sistema nunca foram alteradas sem autorização;
  \item defeitos de concorrência só apareceram sob carga ou em execuções repetidas, o que justificou
        executar a suíte várias vezes seguidas antes de cada entrega;
  \item otimizar a partir de medições, e não de suposições, levou a ganhos grandes com alterações
        pequenas: o custo dominante estava no redesenho da biblioteca gráfica, e não no código Python;
  \item a arquitetura definida na versão 2.0.0 absorveu sem retrabalho estrutural todas as funções
        seguintes (ativação por voz, memória, programas, idiomas, instalador, modelo local, escolha do
        microfone e mudança do banco), cada uma como um adaptador novo atrás de uma porta, o que confirma
        o valor de investir cedo na separação de responsabilidades;
  \item um modelo de linguagem pequeno precisa de limites explícitos: restringir as ferramentas que ele
        pode usar e deixar as ações destrutivas para comandos determinísticos tornou o modelo local útil
        sem expor o usuário a ações indesejadas;
  \item operações sobre os dados do usuário foram desenhadas para nunca destruir informação: o arquivo
        antigo é renomeado, o destino ocupado é recusado e toda cópia é conferida antes de ser usada.
\end{enumerate}
